%% chapter2.tex — Background and Related Work

\chapter{Background and Related Work}
\label{ch:background}

%% -----------------------------------------------------------------------
%% 2.1 Related Work
%% -----------------------------------------------------------------------
\section{Related Work}
\label{sec:related}

This chapter reviews prior work relevant to the thesis topic.
A thorough literature review situates your contribution within the
existing body of knowledge and justifies your methodological choices.

Each paragraph should synthesise related sources rather than
summarising individual papers in isolation~\cite{kopka1999guide}.
Group sources thematically:

\begin{itemize}
  \item Approach A and its variants~\cite{lamport1994latex}.
  \item Approach B and its limitations~\cite{knuth1986texbook}.
  \item Hybrid methods that combine elements of both~\cite{kopka1999guide}.
\end{itemize}

%% -----------------------------------------------------------------------
%% 2.2 Figures
%% -----------------------------------------------------------------------
\section{Including Figures}
\label{sec:figures}

Figures are included with the standard \texttt{figure} environment.
The class automatically adds them to the List of Figures.
Always include a descriptive caption and a \verb|\label| so the figure
can be cross-referenced.

Figure~\ref{fig:example} shows a placeholder rectangle that should be
replaced with your actual figure file (\texttt{.pdf}, \texttt{.png},
or \texttt{.eps}).

\begin{figure}[htbp]
  \centering
  % Replace the \rule with \includegraphics[width=0.8\textwidth]{yourfile}
  \setlength{\fboxsep}{0pt}
  \fbox{\rule{0pt}{6cm}\rule{10cm}{0pt}}
  \caption[Short caption for the List of Figures]{%
    Longer caption that appears below the figure.
    The short caption in brackets is what appears in the List of Figures.
    If only one argument is given it is used in both places.}
  \label{fig:example}
\end{figure}

When a figure is wider than the text column, the \texttt{[width=\textbackslash textwidth]}
option scales it to fit.
Never let a figure extend into the margin.

%% -----------------------------------------------------------------------
%% 2.3 Tables
%% -----------------------------------------------------------------------
\section{Including Tables}
\label{sec:tables}

Tables should use the \texttt{booktabs} package for clean horizontal
rules, as shown in Table~\ref{tab:example}.
Vertical rules are generally discouraged in professional typesetting.

\begin{table}[htbp]
  \centering
  \caption[Short caption for the List of Tables]{%
    A comparison of example methods.  Replace this with your data.}
  \label{tab:example}
  \begin{tabular}{llcc}
    \toprule
    Method      & Reference              & Accuracy (\%) & Runtime (s) \\
    \midrule
    Approach A  & \cite{lamport1994latex}  & 91.2          & 4.3         \\
    Approach B  & \cite{knuth1986texbook}  & 88.7          & 2.1         \\
    Approach C  & \cite{kopka1999guide}    & 93.5          & 8.9         \\
    \midrule
    This work   & ---                    & \textbf{95.1} & 5.7         \\
    \bottomrule
  \end{tabular}
\end{table}

The \texttt{[htbp]} placement specifier asks \LaTeX\ to place the float
\textit{h}ere, at the \textit{t}op of a page, at the \textit{b}ottom,
or on a separate \textit{p}age of floats, in that order of preference.

%% -----------------------------------------------------------------------
%% 2.4 Equations
%% -----------------------------------------------------------------------
\section{Equations}
\label{sec:equations}

Inline mathematics is set with dollar signs: $E = mc^2$.
Displayed equations are set with the \texttt{equation} environment:

\begin{equation}
  \label{eq:gaussian}
  f(x) = \frac{1}{\sigma\sqrt{2\pi}} \exp\!\left(-\frac{(x-\mu)^2}{2\sigma^2}\right).
\end{equation}

Equation~\eqref{eq:gaussian} shows a Gaussian probability density
function with mean $\mu$ and standard deviation $\sigma$.
For aligned multi-line derivations, use the \texttt{align} environment:

\begin{align}
  a + b  &= c, \label{eq:line1} \\
  d \cdot e &= f. \label{eq:line2}
\end{align}

Equations~\eqref{eq:line1} and~\eqref{eq:line2} illustrate aligned
equations.
Equation numbers are set sequentially by the class (or per-chapter
if the \texttt{nonsequential} option is used).
